
\abstract{%

%
% Yujun's version
Offering great potential in robotic manipulation, a capable Vision-Language-Action (VLA) foundation model is expected to faithfully generalize across tasks and platforms while ensuring cost efficiency (\textit{e.g.}, data and GPU hours required for adaptation).
%
To this end, we develop \method with around 20,000 hours of real-world data from 9 popular dual-arm robot configurations.
%
Through a systematic assessment on 4 robotic platforms, each completing 100 tasks with 130 post-training episodes per task, our model achieves clear superiority over competitors, showcasing its \textbf{\textit{strong performance}} and \textbf{\textit{broad generalizability}}.
%
We have also built an \textbf{\textit{efficient}} codebase, which delivers a throughput of $261$ samples per second with an 8-GPU training setup, representing a $1.5\sim2.8\times$ (depending on the relied VLM base model) speedup over existing VLA-oriented codebases.
%
The above features ensure that our model is well-suited for real-world deployment.
%
To advance the field of robot learning, we provide open access to the code, base model, and benchmark data, with a focus on enabling more challenging tasks and promoting sound evaluation standards.
%
%
%
%
%
\iffalse
% Kecheng's version
% Vision-Language-Action (VLA) models enable robots to perform diverse manipulation tasks guided by natural language instructions.
%
A foundational Vision-Language-Action (VLA) model, trained on large-scale and diverse robotic datasets, is crucial for enabling robots to robustly handle real-world variability and rapidly acquire new skills.
%
In this paper, we present a pragmatic cross-embodiment VLA model, named \method, that enables unified control policies across diverse robotic morphologies to accelerate VLA development.
%
Thanks to 20,000 hours data from 7 kinds of real-world robots, we can offer an efficient codebase that supports strong pretrained models to achieve great performance improvements for state-of-the-art VLA models.
%depth?
Using our codebase, a \method model with around 4B parameters can be trained with over 3,000 tokens/sec/GPU throughput on 128 GPUs, showcasing its superior efficiency and scalability.
%
To evaluate the proposed pretrained model, we construct a benchmark, including 100 tasks as well as 3 real-world robots.
%
Code, benchmark data, and pre-trained models will be made publicly available to facilitate the development and evaluation of this field.
\fi
%

}
% \begin{figure}[htbp]
%     \centering
%     \includegraphics[width=0.9\linewidth]{figures/Teaser.pdf}
%     \caption{Caption}
%     \label{fig:data}
% \end{figure}
